\documentclass[11pt]{article}
\usepackage[a4paper,margin=1.05in]{geometry}
\usepackage[T1]{fontenc}
\usepackage{amsmath,amssymb,amsthm,mathtools}
\usepackage{stmaryrd}
\usepackage{mathpartir}
\usepackage{enumitem}
\usepackage{xcolor}
\usepackage[colorlinks=true,linkcolor=blue!55!black,citecolor=blue!55!black,urlcolor=blue!55!black]{hyperref}

%%% Theorem environments
\theoremstyle{plain}
\newtheorem{theorem}{Theorem}[section]
\newtheorem{lemma}[theorem]{Lemma}
\newtheorem{corollary}[theorem]{Corollary}
\newtheorem{proposition}[theorem]{Proposition}
\theoremstyle{definition}
\newtheorem{definition}[theorem]{Definition}
\newtheorem{example}[theorem]{Example}
\newtheorem{assumption}[theorem]{Assumption}
\theoremstyle{remark}
\newtheorem{remark}[theorem]{Remark}

%%% Keywords and syntax
\newcommand{\kw}[1]{\mathsf{#1}}
\newcommand{\kif}{\kw{if}\;}
\newcommand{\kthen}{\;\kw{then}\;}
\newcommand{\kelse}{\;\kw{else}\;}
\newcommand{\nil}{\mathbf{0}}
\newcommand{\ttrue}{\mathsf{tt}}
\newcommand{\ffalse}{\mathsf{ff}}
\newcommand{\seq}{;\;}
\newcommand{\pp}{\mathsf{p}}
\newcommand{\qq}{\mathsf{q}}
\newcommand{\rr}{\mathsf{r}}
\newcommand{\Procs}{\mathcal{P}}
\newcommand{\Var}{\mathit{Var}}
\newcommand{\Val}{\mathit{Val}}
\newcommand{\pn}{\mathrm{pn}}
\newcommand{\eval}[3]{#1 \downarrow_{#2} #3}
\newcommand{\conf}[2]{\langle #1,\, #2\rangle}
\newcommand{\lto}[1]{\xrightarrow{\;#1\;}}
\newcommand{\proj}[2]{\llbracket #1 \rrbracket_{#2}}
\newcommand{\projn}[1]{\llbracket #1 \rrbracket}
\newcommand{\tauat}[1]{\tau\mathord{@}#1}
\newcommand{\cast}[2]{#1 \rightsquigarrow #2}
\newcommand{\upd}[3]{#1[#2 \mapsto #3]}
\newcommand{\rn}[1]{{\normalfont\scshape #1}}
\newcommand{\restr}[2]{#1|_{#2}}
\newcommand{\ifcond}[3]{\kif #1 \kthen #2 \kelse #3}
\newcommand{\sendbr}[4]{#1!#2 \mathbin{\triangleright} (#3,\, #4)}
\newcommand{\recvbr}[3]{#1?(#2,\, #3)}
\newcommand{\pif}[3]{\kif #1 \kthen #2 \kelse #3}
\newcommand{\case}[1]{\medskip\noindent\textbf{Case #1.}\;}
\newcommand{\subcase}[1]{\smallskip\noindent\textit{Subcase #1.}\;}
\newcommand{\Bis}{\mathcal{B}}

%%% Procedures
\newcommand{\Pnames}{\mathcal{X}}
\newcommand{\Defs}{\mathcal{C}}
\newcommand{\NDefs}{\mathcal{D}}
\newcommand{\body}[1]{\Defs(#1)}
\newcommand{\ldef}[2]{\NDefs_{#1}(#2)}
\newcommand{\rt}[3]{#1\mathbin{:}#2.\,#3}
\newcommand{\Loop}{\mathit{Loop}}
\newcommand{\Cif}{C_{\mathrm{if}}}
\newcommand{\CL}{C_{\Loop}}

\title{A Core Choreographic Language with\\Multicast Knowledge of Choice and Recursion}
\author{}
\date{\today}

\begin{document}
\maketitle

\begin{abstract}
  We present a small choreographic programming language in the style of Core
  Choreographies, with local computation, point-to-point communication,
  conditionals and recursive procedures.  The language departs from Core
  Choreographies in how it propagates knowledge of choice.  Instead of
  selections and merging, a conditional guarded by process $\pp$ is
  implemented by having $\pp$ \emph{multicast the value of its guard} to every
  other process that participates in the conditional, including processes
  that participate only through procedures called from its branches, with
  every participant, including $\pp$, selecting its branch in the same step.
  At the network level this is realised by two symmetric atomic constructs,
  \emph{send-and-branch} for the guarding process and
  \emph{receive-and-branch} for the receivers.  Procedure calls, in contrast,
  are entered by each participant independently, and the choreographic
  semantics records a partially entered call with a runtime term.  We give the
  syntax and semantics of choreographies and of network processes, define
  endpoint projection, and prove that projection is correct in both
  directions.  Projection is total (no merging can fail), conditionals need no
  runtime terms, and the correctness theorem is an exact labelled
  bisimulation with no pruning relation.  Its only hypothesis is a
  well-formedness condition on runtime terms, which holds for every program and
  is preserved by execution.
\end{abstract}

\section*{Overview}

Core Choreographies (CC) implements \emph{knowledge of choice} by
\emph{selections}: after a conditional $\ifcond{\pp.e}{C_1}{C_2}$, the process
$\pp$ must inform every other participant of the branch it took with a label
communication of the form $\pp \to \qq[l]$.  Endpoint projection then uses a
partial \emph{merging} operator to combine the two branches of a process that
has not yet learned the choice, and \emph{projectability} becomes a side
condition that may fail.  Moreover, since the projections of the two branches
have been merged, the correspondence between a choreography and its projection
must be stated up to a \emph{pruning} relation that discards the unreachable
alternatives once a choice is known.

This note works out an alternative.  A conditional is realised by a single
synchronous multicast: $\pp$ evaluates its guard to a boolean $v$, sends $v$
to the set $Q$ of all other processes occurring in either branch, and
\emph{in the same step} every participant, $\pp$ included, continues with the
branch selected by $v$.  Processes that occur in neither branch are not
involved at all.  Because every participant learns the choice before executing
any part of a branch, there is nothing to merge and no dead code to prune.

The network-level constructs that implement this are deliberately symmetric.
The guarding process executes a \emph{send-and-branch}
$\sendbr{Q}{e}{B_1}{B_2}$, which evaluates $e$, multicasts the result to $Q$
and continues as $B_1$ or $B_2$ according to that result; each receiver
executes a \emph{receive-and-branch} $\recvbr{\pp}{B_1}{B_2}$, which waits for
a boolean from $\pp$ and continues accordingly.  Since the sender's own
branching is part of the multicast step rather than a separate local action,
there is no intermediate network state in which the announcement has been
made but the sender has not yet committed.  The choreographic semantics
therefore needs no runtime term for a partially executed conditional.

The language includes recursive procedures in the style of CC: a single
global set of definitions $X = C$, with calls in tail position.  Two
aspects of the extension deserve comment.  First, the receivers of a
conditional are now all processes that may take part in either branch,
\emph{including through procedures called from it}.  This set is computed
statically, as a least fixed point over the procedure definitions, and the
multicast mechanism is otherwise unchanged.  Second, a procedure call is not
atomic at the network level: each process unfolds its own local version of
the procedure, independently of the others.  We make every such unfolding a
visible step, both in networks and in choreographies, and record a call that
has been entered by some but not all of its participants with a runtime term
$\rt{P}{X}{C}$, in which $P$ is the set of processes still to enter and $C$
is the current state of the body.  This is the only runtime term in the
language.  The price is a well-formedness invariant on runtime terms, which
holds trivially for programs and is preserved by execution; under it, the
correspondence between choreographies and networks is a step-for-step
labelled bisimulation, with no pruning and no unfolding up to a structural
congruence.

The remainder is organised as follows: \S\ref{sec:chor} gives the syntax and
semantics of choreographies, \S\ref{sec:net} the syntax and semantics of
network processes, \S\ref{sec:epp} endpoint projection, and
\S\ref{sec:correct} the proof of correctness in both directions.

\section{Choreographies}
\label{sec:chor}

\subsection{Preliminaries}

We fix a finite set $\Procs$ of \emph{process names}, ranged over by
$\pp,\qq,\rr$; sets of process names are ranged over by $P$ and $Q$.  We fix a
finite set $\Pnames$ of \emph{procedure names}, ranged over by $X,Y$.  We also
fix a set $\Var$ of \emph{variables}, ranged over by $x,y,z$, and a set $\Val$
of \emph{values}, ranged over by $v,w$, containing the booleans $\ttrue$ and
$\ffalse$.

Local \emph{expressions} $e$ are a parameter of the language.  A \emph{local
store} is a total function $s : \Var \to \Val$.  We write
$\eval{e}{s}{v}$ when $e$ evaluates to $v$ in store $s$.

\begin{assumption}\label{ass:eval}
  Evaluation is total and deterministic: for every $e$ and $s$ there is exactly
  one $v$ with $\eval{e}{s}{v}$.  Furthermore, every expression used as the
  guard of a conditional evaluates to a boolean in every store.
\end{assumption}

\noindent
The second half of the assumption is a typing concern that is orthogonal to the
present development; any standard local type system enforces it.  A (global)
\emph{store} is a function $\sigma : \Procs \to (\Var \to \Val)$; we write
$\sigma_\pp$ for $\sigma(\pp)$ and $\upd{\sigma}{\pp.x}{v}$ for the store that
agrees with $\sigma$ except that $\upd{\sigma}{\pp.x}{v}_\pp(x) = v$.

\subsection{Syntax}

\begin{definition}[Choreographies]\label{def:chor-syntax}
  Actions $\eta$ and choreographies $C$ are given by the grammar
  \[
  \begin{array}{r@{\;\;}c@{\;\;}l@{\qquad}l}
    \eta & ::= & \pp.x := e & \text{(local computation)}\\
         & \mid & \pp.e \to \qq.x & \text{(communication, $\pp \neq \qq$)}\\[1ex]
    C & ::= & \nil & \text{(terminated)}\\
      & \mid & \eta \seq C & \text{(action)}\\
      & \mid & \ifcond{\pp.e}{C_1}{C_2} & \text{(conditional)}\\
      & \mid & X & \text{(procedure call)}\\
      & \mid & \rt{P}{X}{C} & \text{(runtime term, $\emptyset \neq P \subseteq \Procs$)}
  \end{array}
  \]
  For a conditional with branches $C_1, C_2$ we write $C_\ttrue = C_1$ and
  $C_\ffalse = C_2$, so that $C_v$ is the branch selected by the boolean $v$.
  A choreography is a \emph{source choreography} if it contains no runtime
  term.
\end{definition}

\begin{definition}[Procedure definitions]\label{def:defs}
  A set of \emph{procedure definitions} $\Defs$ assigns to every procedure
  name $X\in\Pnames$ a source choreography $\body{X}$, the \emph{body} of $X$.
  We write $X = C$ when $\body{X} = C$.  Throughout, we fix one set of
  procedure definitions $\Defs$.
\end{definition}

The action $\pp.x := e$ has $\pp$ evaluate $e$ and store the result in its
local variable $x$.  The action $\pp.e \to \qq.x$ has $\pp$ evaluate $e$ and
send the result to $\qq$, which stores it in $x$.  In
$\ifcond{\pp.e}{C_1}{C_2}$ the process $\pp$ evaluates its guard $e$ and the
choreography continues as $C_1$ or $C_2$ accordingly.  A call $X$ continues
as the body $\body{X}$.  Bodies may call any procedure, including the one
being defined, so procedures may be recursive and mutually recursive.  As in
CC, conditionals and calls occur only in tail position: there is no
sequential composition after either, and a continuation shared by both
branches of a conditional is written in each of them.  In particular, a loop
is a procedure whose body calls the procedure again in some branches and
continues with whatever follows the loop in the others.  Unlike CC, there are
no selections: knowledge of choice is built into the semantics of the
conditional itself.

The runtime term $\rt{P}{X}{C}$ is not written by programmers.  It arises
during execution and records a call to $X$ that some, but not all, of its
participants have entered: the body of $X$ has progressed to $C$, and the
processes in $P$ have yet to enter the call.  Its role is explained together
with the semantics in \S\ref{sec:chor-sem}.  So that the rules can be stated
uniformly, we adopt the following convention: when $P = \emptyset$, the
notation $\rt{P}{X}{C}$ stands for $C$ itself.  Once every participant has
entered a call, there is nothing left to record.

Procedures here take no parameters, as in CC.  Procedures parameterised on
process names, as in Montesi's textbook presentation, would additionally need
substitution of process names in choreographies and behaviours, and
projection of procedure bodies onto parameter positions; we do not treat them.

The process names occurring in a choreography now depend on the procedures it
calls, and since procedures may be recursive these are given by a fixed point.

\begin{definition}[Process names]\label{def:pn}
  For a function $\rho : \Pnames \to 2^{\Procs}$, the set $\pn_\rho(\cdot)$ of
  process names occurring in an action or choreography is defined by
  \[
  \begin{array}{l@{\quad}l}
    \pn(\pp.x := e) = \{\pp\} &
    \pn(\pp.e \to \qq.x) = \{\pp,\qq\}\\[0.5ex]
    \pn_\rho(\nil) = \emptyset &
    \pn_\rho(\eta \seq C) = \pn(\eta) \cup \pn_\rho(C)\\[0.5ex]
    \pn_\rho(X) = \rho(X) &
    \pn_\rho(\rt{P}{X}{C}) = P \cup \pn_\rho(C)\\[0.5ex]
    \multicolumn{2}{l}{
    \pn_\rho(\ifcond{\pp.e}{C_1}{C_2}) = \{\pp\}\cup\pn_\rho(C_1)\cup\pn_\rho(C_2).}
  \end{array}
  \]
  Let $\pi : \Pnames \to 2^{\Procs}$ be the least function, ordered pointwise
  by inclusion, such that $\pi(X) = \pn_\pi(\body{X})$ for every
  $X\in\Pnames$.  We write $\pn(C)$ for $\pn_\pi(C)$; in particular
  $\pn(X) = \pi(X)$, the set of \emph{participants} of $X$.  The
  \emph{receivers} of a conditional at $\pp$ with branches $C_1,C_2$ are the
  processes $(\pn(C_1)\cup\pn(C_2)) \setminus \{\pp\}$.
\end{definition}

The function $\pi$ exists because $\pn_\rho(C)$ is monotone in $\rho$, so
the map $\rho \mapsto (X \mapsto \pn_\rho(\body{X}))$ is monotone on the
finite lattice of functions $\Pnames\to 2^{\Procs}$.  It is computed by
iterating this map from the function that is everywhere $\emptyset$, which
stabilises after finitely many steps.  Concretely, $\pn(X)$ is the set of
processes that occur in an action, or as the guard of a conditional, in the
body of $X$ or of some procedure reachable from $X$ through calls.

\begin{lemma}[Participants of a procedure]\label{lem:pn-body}
  $\pn(\body{X}) = \pn(X)$ for every $X\in\Pnames$.
\end{lemma}
\begin{proof}
  $\pn(\body{X}) = \pn_\pi(\body{X}) = \pi(X) = \pn(X)$, since $\pi$ is a
  fixed point.
\end{proof}

The proofs below use only this fixed-point property; taking the \emph{least}
fixed point keeps the sets of receivers, and hence the number of messages, as
small as possible.  Note that the receivers of a conditional are all
processes other than $\pp$ that occur in \emph{either} branch, even if they
occur in only one of them, or only in a procedure called from one of them.
The receiver set may be empty.

\begin{example}\label{ex:chor}
  Let $\Procs = \{\pp,\qq,\rr\}$ and $\Pnames = \{\Loop\}$ with
  \[
  \begin{array}{l@{\;}l}
    \Loop \;=\; \pp.x := 2\cdot x \seq
    & \kif \pp.(x > 10) \kthen
       \big(\rr.n := n+1 \seq \pp.x \to \qq.y \seq \qq.y \to \rr.z \seq \nil\big)\\
    & \kw{else}\; \big(\qq.k := k+1 \seq \Loop\big).
  \end{array}
  \]
  Process $\pp$ doubles $x$ until it exceeds $10$, while $\qq$ counts the
  number of rounds that did not end the loop.  When the loop ends, $\rr$
  increments a counter and $\pp$'s value is forwarded to $\rr$ via $\qq$.  We
  take the call $\Loop$ as the program.  Starting from the everywhere-empty
  function, one iteration gives $\{\pp,\qq,\rr\}$ for $\Loop$ and a second
  iteration gives the same set, so $\pn(\Loop) = \{\pp,\qq,\rr\}$.  The
  receivers of the conditional are $\{\qq,\rr\}$; note that $\rr$ occurs only
  in the first branch and in the call in the second.  We abbreviate
  \[
    \CL = \body{\Loop} = \pp.x := 2\cdot x \seq \Cif,
    \qquad
    \Cif = \ifcond{\pp.(x>10)}{C_1}{C_2},
  \]
  with $C_1 = \rr.n := n+1 \seq \pp.x \to \qq.y \seq \qq.y \to \rr.z \seq \nil$
  and $C_2 = \qq.k := k+1 \seq \Loop$.
\end{example}

\subsection{Semantics}
\label{sec:chor-sem}

The semantics is a labelled transition system over \emph{configurations}
$\conf{C}{\sigma}$.  Transition labels record which processes take part in a
step and what kind of step it is.

\begin{definition}[Labels]\label{def:labels}
  Labels $\mu$ are given by
  \[
    \mu ::= \tauat{\pp} \;\mid\; \pp \to \qq \;\mid\; \cast{\pp}{Q}
  \]
  denoting, respectively, an internal step at $\pp$, a communication from
  $\pp$ to $\qq$, and a multicast from $\pp$ to the set $Q$ (with $\pp\notin Q$).
  The processes involved in a label are
  $\pn(\tauat{\pp}) = \{\pp\}$, $\pn(\pp\to\qq) = \{\pp,\qq\}$ and
  $\pn(\cast{\pp}{Q}) = \{\pp\}\cup Q$.
\end{definition}

\begin{definition}[Choreographic semantics]\label{def:chor-sem}
  The relation $\conf{C}{\sigma} \lto{\mu} \conf{C'}{\sigma'}$ is the least
  relation closed under the rules in Figure~\ref{fig:chor-sem}.
\end{definition}

\begin{figure}[t]
\begin{mathpar}
\inferrule*[lab=\rn{C|Local}]
  {\eval{e}{\sigma_\pp}{v}}
  {\conf{\pp.x := e \seq C}{\sigma} \lto{\tauat{\pp}} \conf{C}{\upd{\sigma}{\pp.x}{v}}}
\and
\inferrule*[lab=\rn{C|Com}]
  {\eval{e}{\sigma_\pp}{v}}
  {\conf{\pp.e \to \qq.x \seq C}{\sigma} \lto{\pp\to\qq} \conf{C}{\upd{\sigma}{\qq.x}{v}}}
\\
\inferrule*[lab=\rn{C|Cast}]
  {\eval{e}{\sigma_\pp}{v} \\ Q = (\pn(C_1)\cup\pn(C_2))\setminus\{\pp\}}
  {\conf{\ifcond{\pp.e}{C_1}{C_2}}{\sigma} \lto{\cast{\pp}{Q}} \conf{C_v}{\sigma}}
\and
\inferrule*[lab=\rn{C|Delay}]
  {\conf{C}{\sigma} \lto{\mu} \conf{C'}{\sigma'} \\ \pn(\eta)\cap\pn(\mu) = \emptyset}
  {\conf{\eta \seq C}{\sigma} \lto{\mu} \conf{\eta \seq C'}{\sigma'}}
\\
\inferrule*[lab=\rn{C|Call}]
  {\pp\in\pn(X)}
  {\conf{X}{\sigma} \lto{\tauat{\pp}} \conf{\rt{(\pn(X)\setminus\{\pp\})}{X}{\body{X}}}{\sigma}}
\and
\inferrule*[lab=\rn{C|Enter}]
  {\pp\in P}
  {\conf{\rt{P}{X}{C}}{\sigma} \lto{\tauat{\pp}} \conf{\rt{(P\setminus\{\pp\})}{X}{C}}{\sigma}}
\\
\inferrule*[lab=\rn{C|Inside}]
  {\conf{C}{\sigma} \lto{\mu} \conf{C'}{\sigma'} \\ P\cap\pn(\mu) = \emptyset}
  {\conf{\rt{P}{X}{C}}{\sigma} \lto{\mu} \conf{\rt{P}{X}{C'}}{\sigma'}}
\end{mathpar}
\caption{Semantics of choreographies.  In \rn{C|Call} and \rn{C|Enter}, a
target of the form $\rt{\emptyset}{X}{C}$ stands for $C$.}
\label{fig:chor-sem}
\end{figure}

Rules \rn{C|Local} and \rn{C|Com} are as in CC.  Rule \rn{C|Cast} executes a
conditional in a single step: the guard is evaluated at $\pp$, the resulting
boolean $v$ is announced to the receivers, and the choreography continues as
the selected branch $C_v$.  The label $\cast{\pp}{Q}$ records the announcement
and its recipients; the store is unchanged.

Rule \rn{C|Delay} implements \emph{out-of-order execution}, which is needed
because a network of independent processes may execute actions of different
processes in any order.  It lets a step of the continuation $C$ happen before
the action $\eta$, provided the processes involved are disjoint.

The remaining three rules execute procedure calls, one participant at a
time.  Rule \rn{C|Call} lets a participant $\pp$ of $X$ enter the call: the
body of $X$ becomes active, and every other participant still has to enter.
Rule \rn{C|Enter} lets a further participant enter; when the last one does,
the runtime term disappears by our convention.  Rule \rn{C|Inside} lets the
processes that have already entered proceed with the body, as long as the
step involves no process that has not.  It plays the role of \rn{C|Delay} for
runtime terms, with the set $P$ in place of the action $\eta$.  Entering a
call changes no store and is labelled as an internal step of the entering
process.

\begin{remark}[No delay through a conditional]\label{rem:no-delay-cond}
  CC has a rule allowing a step to be delayed past a conditional when both
  branches can perform that same step, which is exactly the situation in
  which merging is needed.  There is no such rule here, and none is needed.
  Every receiver is waiting for $\pp$'s announcement and cannot act before
  it, while any process that is not a receiver does not occur in either branch
  and so has nothing to do.  The only step a configuration
  $\conf{\ifcond{\pp.e}{C_1}{C_2}}{\sigma}$ can take is \rn{C|Cast}.  This is
  confirmed by the network semantics in \S\ref{sec:net} and by the soundness
  proof in \S\ref{sec:correct}.  When a conditional appears inside a runtime
  term $\rt{P}{X}{C}$, rule \rn{C|Inside} additionally prevents the multicast
  until every receiver has entered the call, since a receiver in $P$ is not
  yet waiting for the announcement.
\end{remark}

\begin{remark}[Calls are entered one process at a time]\label{rem:entry}
  In the network semantics of \S\ref{sec:net}, each process unfolds its own
  local version of a procedure, without synchronising with the other
  participants.  If a choreographic call unfolded the body for all
  participants at once, as it does when calls are unfolded by a structural
  congruence, that single step would correspond to several network steps, or
  to none, and the correspondence between choreographies and networks could
  only be stated up to unfolding.  Here every unfolding is a labelled step at
  both levels.  The runtime term is what makes this possible: participants
  may enter a call at different times, and those that have entered may make
  progress in the body in the meantime.  It is the only runtime term in the
  language, since a conditional is still resolved in a single step.
\end{remark}

\begin{example}\label{ex:run}
  Continuing Example~\ref{ex:chor}, let $\sigma$ be a store with
  $\sigma_\pp(x) = 3$, $\sigma_\qq(k) = 0$ and $\sigma_\rr(n) = 0$.  One
  execution of the program $\Loop$ begins
  \begin{align*}
    \conf{\Loop}{\sigma}
      &\lto{\tauat{\pp}} \conf{\rt{\{\qq,\rr\}}{\Loop}{\CL}}{\sigma}
      && \text{by \rn{C|Call}}\\
      &\lto{\tauat{\pp}} \conf{\rt{\{\qq,\rr\}}{\Loop}{\Cif}}{\sigma_1}
      && \text{by \rn{C|Inside}, \rn{C|Local}}\\
      &\lto{\tauat{\qq}} \conf{\rt{\{\rr\}}{\Loop}{\Cif}}{\sigma_1}
      && \text{by \rn{C|Enter}}\\
      &\lto{\tauat{\rr}} \conf{\Cif}{\sigma_1}
      && \text{by \rn{C|Enter}}\\
      &\lto{\cast{\pp}{\{\qq,\rr\}}} \conf{\qq.k := k+1 \seq \Loop}{\sigma_1}
      && \text{by \rn{C|Cast}}\\
      &\lto{\tauat{\pp}} \conf{\qq.k := k+1 \seq \rt{\{\qq,\rr\}}{\Loop}{\CL}}{\sigma_1}
      && \text{by \rn{C|Delay}, \rn{C|Call}}\\
      &\lto{\tauat{\pp}} \conf{\qq.k := k+1 \seq \rt{\{\qq,\rr\}}{\Loop}{\Cif}}{\sigma_2}
      && \text{by \rn{C|Delay}, \rn{C|Inside}, \rn{C|Local}}\\
      &\lto{\tauat{\rr}} \conf{\qq.k := k+1 \seq \rt{\{\qq\}}{\Loop}{\Cif}}{\sigma_2}
      && \text{by \rn{C|Delay}, \rn{C|Enter}}\\
      &\lto{\tauat{\qq}} \conf{\rt{\{\qq\}}{\Loop}{\Cif}}{\sigma_3}
      && \text{by \rn{C|Local}}\\
      &\lto{\tauat{\qq}} \conf{\Cif}{\sigma_3}
      && \text{by \rn{C|Enter}}\\
      &\lto{\cast{\pp}{\{\qq,\rr\}}} \conf{C_1}{\sigma_3}
      && \text{by \rn{C|Cast}}
  \end{align*}
  where $\sigma_1 = \upd{\sigma}{\pp.x}{6}$,
  $\sigma_2 = \upd{\sigma_1}{\pp.x}{12}$ and
  $\sigma_3 = \upd{\sigma_2}{\qq.k}{1}$.  The first multicast announces
  $\ffalse$, since $6 \not> 10$, and the second announces $\ttrue$.  From
  $\conf{C_1}{\sigma_3}$ three further steps, by \rn{C|Local}, \rn{C|Com}
  and \rn{C|Com}, reach $\nil$ with $\rr$ holding $z = 12$.

  In the second step $\pp$ doubles $x$ before $\qq$ and $\rr$ have entered
  the call.  After the first multicast, $\pp$ re-enters $\Loop$ and doubles
  $x$ again while $\qq$ has not yet incremented $k$ for the previous round,
  and $\rr$ enters as well.  At that point the second multicast is not
  enabled: $\qq$ has not entered the call, and both \rn{C|Delay} and
  \rn{C|Inside} forbid a step involving $\qq$.  It becomes enabled once $\qq$
  has caught up.  Many other interleavings are possible; for instance $\qq$
  and $\rr$ may enter the first call in either order, and $\rr$ may enter the
  second call before or after $\pp$ doubles $x$.
\end{example}

\begin{remark}[Nesting of runtime terms]\label{rem:nesting}
  A process may fall arbitrarily far behind the others, and runtime terms may
  then nest.  With $X = \pp.e \to \qq.x \seq \rr.y := 1 \seq X$, processes
  $\pp$ and $\qq$ can iterate the loop while $\rr$ has not entered even the
  first call.  After two rounds the choreography is
  \[
    \rt{\{\rr\}}{X}{\big(\rr.y := 1 \seq
      \rt{\{\rr\}}{X}{(\rr.y := 1 \seq X)}\big)},
  \]
  which records the actions that $\rr$ still owes.  This growth is a
  consequence of out-of-order execution and arises equally when calls are
  unfolded by a structural congruence.  The network state, by contrast, stays
  small: $\rr$ is still at its call to $X$.
\end{remark}

\begin{lemma}[Store locality]\label{lem:chor-locality}
  If $\conf{C}{\sigma} \lto{\mu} \conf{C'}{\sigma'}$ then
  $\sigma'_\rr = \sigma_\rr$ for every $\rr \notin \pn(\mu)$.
\end{lemma}
\begin{proof}
  By induction on the derivation.  \rn{C|Local} changes the store only at
  $\pp \in \pn(\tauat\pp)$, and \rn{C|Com} only at $\qq \in \pn(\pp\to\qq)$.
  \rn{C|Cast}, \rn{C|Call} and \rn{C|Enter} leave the store unchanged.
  \rn{C|Delay} and \rn{C|Inside} preserve the label and the store change of
  their premise, so the claim follows from the induction hypothesis.
\end{proof}

\section{Network Processes}
\label{sec:net}

\subsection{Syntax}

\begin{definition}[Processes and networks]\label{def:net-syntax}
  Behaviours $B$ are given by the grammar
  \[
  \begin{array}{r@{\;\;}c@{\;\;}l@{\qquad}l}
    B & ::= & \nil & \text{(terminated)}\\
      & \mid & x := e \seq B & \text{(local computation)}\\
      & \mid & \qq!e \seq B & \text{(send to $\qq$)}\\
      & \mid & \pp?x \seq B & \text{(receive from $\pp$ into $x$)}\\
      & \mid & \sendbr{Q}{e}{B_1}{B_2} & \text{(send-and-branch: multicast $e$ to $Q$, branch on it)}\\
      & \mid & \recvbr{\pp}{B_1}{B_2} & \text{(receive-and-branch: receive a boolean from $\pp$, branch on it)}\\
      & \mid & X & \text{(procedure call)}
  \end{array}
  \]
  A set of \emph{local procedure definitions} $\NDefs$ assigns to every
  process $\rr\in\Procs$ and procedure name $X\in\Pnames$ a behaviour
  $\ldef{\rr}{X}$, the body of $\rr$'s local procedure $X$.
  A \emph{network} $N$ is a total function from $\Procs$ to behaviours.  We
  write $\pp[B]$ for the network fragment mapping $\pp$ to $B$, and
  $\pp[B] \mid N$ for the network that maps $\pp$ to $B$ and agrees with $N$
  on all other processes (where $N$ is then understood as a function on
  $\Procs\setminus\{\pp\}$).  Parallel composition $\mid$ is therefore
  commutative and associative.  For a set $S\subseteq\Procs$ we write
  $\restr{N}{S}$ for the restriction of $N$ to $S$, and we say that $N$ and
  $M$ \emph{agree on $S$} if $\restr{N}{S} = \restr{M}{S}$.
\end{definition}

The behaviours $x := e\seq B$, $\qq!e\seq B$ and $\pp?x\seq B$ are the usual
local assignment, send and receive.  The send-and-branch and
receive-and-branch are the network-level counterparts of a choreographic
conditional and are symmetric to each other.  The \emph{send-and-branch}
$\sendbr{Q}{e}{B_1}{B_2}$ evaluates $e$ to a boolean, multicasts it to every
process in $Q$, and continues as $B_1$ if the boolean is $\ttrue$ and as $B_2$
if it is $\ffalse$.  The \emph{receive-and-branch} $\recvbr{\pp}{B_1}{B_2}$
waits for a boolean from $\pp$ and continues as $B_1$ or $B_2$ accordingly.
There is no separate local conditional: a send-and-branch with
$Q = \emptyset$ evaluates its guard and branches without communicating, and
we may write $\pif{e}{B_1}{B_2}$ as an abbreviation for
$\sendbr{\emptyset}{e}{B_1}{B_2}$.  A call $X$ at process $\rr$ continues as
$\rr$'s own body $\ldef{\rr}{X}$; each process has its own version of every
procedure.  Both branching constructs and calls occur only in tail position.
A network is \emph{terminated} if it maps every process to $\nil$.

\subsection{Semantics}

\begin{definition}[Network semantics]\label{def:net-sem}
  Fix a set of local procedure definitions $\NDefs$.  Network configurations
  are pairs $\conf{N}{\sigma}$ of a network and a store.  The relation
  $\conf{N}{\sigma}\lto{\mu}\conf{N'}{\sigma'}$, with labels as in
  Definition~\ref{def:labels}, is the least relation closed under the rules
  in Figure~\ref{fig:net-sem}.  In rule \rn{N|Com} the processes $\pp$ and
  $\qq$ are distinct; in rule \rn{N|Cast} the processes
  $\pp,\qq_1,\ldots,\qq_n$ are pairwise distinct and $n\geq 0$.
\end{definition}

\begin{figure}[t]
\begin{mathpar}
\inferrule*[lab=\rn{N|Local}]
  {\eval{e}{\sigma_\pp}{v}}
  {\conf{\pp[x := e \seq B] \mid N}{\sigma} \lto{\tauat\pp} \conf{\pp[B]\mid N}{\upd{\sigma}{\pp.x}{v}}}
\and
\inferrule*[lab=\rn{N|Com}]
  {\eval{e}{\sigma_\pp}{v}}
  {\conf{\pp[\qq!e \seq B] \mid \qq[\pp?x \seq B'] \mid N}{\sigma} \lto{\pp\to\qq} \conf{\pp[B]\mid\qq[B']\mid N}{\upd{\sigma}{\qq.x}{v}}}
\\
\inferrule*[lab=\rn{N|Cast}]
  {\eval{e}{\sigma_\pp}{v} \\ Q = \{\qq_1,\ldots,\qq_n\}}
  {\conf{\pp[\sendbr{Q}{e}{B^1}{B^2}] \mid \textstyle\prod_{i=1}^{n}\qq_i[\recvbr{\pp}{B_i^1}{B_i^2}] \mid N}{\sigma}
   \lto{\cast{\pp}{Q}}
   \conf{\pp[B^v] \mid \textstyle\prod_{i=1}^{n}\qq_i[B_i^v] \mid N}{\sigma}}
\\
\inferrule*[lab=\rn{N|Call}]
  { }
  {\conf{\pp[X] \mid N}{\sigma} \lto{\tauat\pp} \conf{\pp[\ldef{\pp}{X}] \mid N}{\sigma}}
\end{mathpar}
\caption{Semantics of networks.  In \rn{N|Cast}, $B^v$ denotes $B^1$ if $v=\ttrue$ and $B^2$ if $v=\ffalse$, and likewise for $B_i^v$.}
\label{fig:net-sem}
\end{figure}

Communication is synchronous.  Rule \rn{N|Cast} is a synchronous
one-to-many rendezvous: in a single step the sender evaluates its guard,
every receiver in $Q$ receives the boolean, and the sender and all receivers
select their continuations.  With $Q=\emptyset$ the rule degenerates to a
local conditional at $\pp$, still labelled $\cast{\pp}{\emptyset}$.  Rule
\rn{N|Call} unfolds a call locally, without synchronising with any other
process; it is labelled as an internal step.  No network counterpart of the
choreographic runtime term is needed: a call that has been entered by some
of its participants is simply a network in which some processes are still
at the call and the others are executing their local bodies.  No structural
rule for parallel composition is needed because networks are functions and
each rule carries the untouched part of the network as $N$.

\section{Endpoint Projection}
\label{sec:epp}

\begin{definition}[Endpoint projection]\label{def:epp}
  The projection $\proj{C}{\rr}$ of a choreography $C$ onto a process
  $\rr\in\Procs$ is defined by structural recursion on $C$:
  \[
  \begin{array}{l@{\;\;}l@{\quad}l}
    \proj{\nil}{\rr} & = \nil \\[1ex]
    \proj{\pp.x := e \seq C}{\rr} & = x := e \seq \proj{C}{\rr} & \text{if } \rr = \pp\\
                                  & = \proj{C}{\rr} & \text{otherwise}\\[1ex]
    \proj{\pp.e \to \qq.x \seq C}{\rr} & = \qq!e \seq \proj{C}{\rr} & \text{if } \rr = \pp\\
                                       & = \pp?x \seq \proj{C}{\rr} & \text{if } \rr = \qq\\
                                       & = \proj{C}{\rr} & \text{otherwise}\\[1ex]
    \proj{\ifcond{\pp.e}{C_1}{C_2}}{\rr} & = \sendbr{Q}{e}{\proj{C_1}{\pp}}{\proj{C_2}{\pp}} & \text{if } \rr = \pp\\
                                         & = \recvbr{\pp}{\proj{C_1}{\rr}}{\proj{C_2}{\rr}} & \text{if } \rr \in Q\\
                                         & = \nil & \text{otherwise}\\[1ex]
    \proj{X}{\rr} & = X & \text{if } \rr \in \pn(X)\\
                  & = \nil & \text{otherwise}\\[1ex]
    \proj{\rt{P}{X}{C}}{\rr} & = X & \text{if } \rr \in P\\
                             & = \proj{C}{\rr} & \text{otherwise}
  \end{array}
  \]
  where, in the conditional case, $Q = (\pn(C_1)\cup\pn(C_2))\setminus\{\pp\}$
  is the set of receivers.  The projection of $C$ is the network
  $\projn{C}$ with $\projn{C}(\rr) = \proj{C}{\rr}$ for every $\rr\in\Procs$,
  and the projection of a configuration is
  $\projn{C,\sigma} = \conf{\projn{C}}{\sigma}$.  The projection of the
  procedure definitions $\Defs$ is the set of local procedure definitions
  $\projn{\Defs}$ given by $\projn{\Defs}_\rr(X) = \proj{\body{X}}{\rr}$ for
  every $\rr\in\Procs$ and $X\in\Pnames$.
\end{definition}

The cases for $\nil$ and for actions are standard.  A conditional at $\pp$
projects onto $\pp$ as a send-and-branch whose target set is the receivers and
whose continuations are the projections of the two branches onto $\pp$.  It
projects onto a receiver $\rr\in Q$ as a receive-and-branch from $\pp$ whose
two continuations are the projections of the two branches onto $\rr$; if $\rr$
does not occur in one of the branches, the corresponding continuation is
$\nil$.  A process that is neither $\pp$ nor a receiver projects to $\nil$,
because it occurs in neither branch.

A call projects onto each participant of the procedure as a call to that
participant's local version of the procedure, and onto every other process as
$\nil$.  Calls are not unfolded by projection: the bodies are projected
separately, once, to obtain $\projn{\Defs}$.  A runtime term $\rt{P}{X}{C}$
projects onto a process that has not yet entered the call as the call $X$,
and onto every other process as the projection of the current state $C$ of
the body.

\begin{remark}[Totality]\label{rem:total}
  Projection is a total function, both on choreographies and on procedure
  definitions: every choreography is projectable.  In CC, the projection of a
  conditional onto a process that is not the guard requires merging the
  projections of the two branches, an operation that is only defined when the
  two behaviours agree up to the selections that distinguish them; programs
  in which a process behaves differently in the two branches without first
  being informed by a selection are rejected.  Here the announcement of the
  guard is built into the semantics of the conditional and is directed to
  \emph{every} process that occurs in either branch, including processes
  that occur only through a procedure call, so the situation that merging is
  designed to handle never arises.  The price is that every participant of a
  conditional always receives a message, even one whose behaviour is
  identical in both branches.
\end{remark}

\begin{example}\label{ex:epp}
  The projection of the program $\Loop$ from Example~\ref{ex:chor} is the
  network $\pp[\Loop] \mid \qq[\Loop] \mid \rr[\Loop]$, and the projected
  local procedure definitions are
  \[
  \begin{array}{l@{\;=\;}l}
    \projn{\Defs}_\pp(\Loop) & x := 2\cdot x \seq \sendbr{\{\qq,\rr\}}{(x > 10)}{\qq!x \seq \nil}{\Loop}\\
    \projn{\Defs}_\qq(\Loop) & \recvbr{\pp}{\pp?y \seq \rr!y \seq \nil}{k := k+1 \seq \Loop}\\
    \projn{\Defs}_\rr(\Loop) & \recvbr{\pp}{n := n+1 \seq \qq?z \seq \nil}{\Loop}.
  \end{array}
  \]
  Process $\rr$ acts only in the round that ends the loop, but as a receiver
  of the conditional it receives the guard in every round, and in a round
  that does not end the loop its only action is to call $\Loop$ again.  In CC
  this cost would not be avoided: $\rr$'s behaviours in the two branches
  differ, so a selection informing $\rr$ would be needed in every round as
  well.

  Reading the execution in Example~\ref{ex:run} through the projection,
  after the first multicast $\pp$ and $\rr$ are at the call $\Loop$ and $\qq$
  is at $k := k+1 \seq \Loop$, which is exactly the projection of $C_2$.
  After the next three steps, in which $\pp$ unfolds its call and doubles $x$
  and $\rr$ unfolds its call, $\pp$ is at a send-and-branch to $\{\qq,\rr\}$
  and $\rr$ at a receive-and-branch, but $\qq$ is still at
  $k := k+1 \seq \Loop$, so \rn{N|Cast} does not apply.  This network is
  exactly the projection of
  $\qq.k := k+1 \seq \rt{\{\qq\}}{\Loop}{\Cif}$.
\end{example}

\section{Correctness of Endpoint Projection}
\label{sec:correct}

Throughout this section, the network semantics is taken with respect to the
projected definitions $\NDefs = \projn{\Defs}$, so that
$\ldef{\rr}{X} = \proj{\body{X}}{\rr}$.  We prove that a well-formed
choreography and its projection have exactly the same labelled transitions,
and that the projection of the target of a choreographic step is the target
of the corresponding network step.  Well-formedness
(\S\ref{sec:wf}) constrains runtime terms only; it holds for every source
choreography and is preserved by execution.

\subsection{Auxiliary lemmas}

\begin{lemma}[Network locality and framing]\label{lem:frame}
  Suppose $\conf{N}{\sigma}\lto{\mu}\conf{N'}{\sigma'}$.
  \begin{enumerate}[label=(\roman*),nosep]
  \item\label{it:frame-loc} $N'(\rr) = N(\rr)$ and $\sigma'_\rr = \sigma_\rr$ for every $\rr\notin\pn(\mu)$.
  \item\label{it:frame-frame} For every network $M$ that agrees with $N$ on
    $\pn(\mu)$, $\conf{M}{\sigma}\lto{\mu}\conf{M'}{\sigma'}$, where $M'$
    agrees with $N'$ on $\pn(\mu)$ and with $M$ on $\Procs\setminus\pn(\mu)$.
  \end{enumerate}
\end{lemma}
\begin{proof}
  By inspection of Figure~\ref{fig:net-sem}.  In each rule, the processes
  displayed explicitly on the left-hand side are exactly the processes in
  $\pn(\mu)$: $\pp$ for \rn{N|Local} and \rn{N|Call}, $\pp$ and $\qq$ for
  \rn{N|Com}, and $\pp,\qq_1,\ldots,\qq_n$ for \rn{N|Cast}.  The rest of the
  network is passed through as $N$ unchanged, and the store is modified only
  at a process in $\pn(\mu)$, which gives~\ref{it:frame-loc}.
  For~\ref{it:frame-frame}, the premises of each rule mention only the
  behaviours of the processes in $\pn(\mu)$, the store $\sigma$ and the fixed
  definitions $\NDefs$; if $M$ agrees with $N$ on $\pn(\mu)$ then the same
  rule instance applies to $\conf{M}{\sigma}$, with the remainder
  $\restr{M}{\Procs\setminus\pn(\mu)}$ in place of
  $\restr{N}{\Procs\setminus\pn(\mu)}$, yielding the same label, the same
  behaviours at $\pn(\mu)$ and the same store $\sigma'$.
\end{proof}

\begin{lemma}[Inversion for networks]\label{lem:inversion}
  Suppose $\conf{N}{\sigma}\lto{\mu}\conf{N'}{\sigma'}$ and $\rr\in\pn(\mu)$.
  Then $N(\rr)\neq\nil$, and exactly one of the following holds.
  \begin{enumerate}[label=(\alph*),nosep]
  \item\label{it:inv-local} $N(\rr) = x := e\seq B$; then the rule is \rn{N|Local}, $\mu = \tauat\rr$, $N'(\rr) = B$, and $\sigma' = \upd{\sigma}{\rr.x}{v}$ where $\eval{e}{\sigma_\rr}{v}$.
  \item\label{it:inv-send} $N(\rr) = \qq!e\seq B$; then the rule is \rn{N|Com}, $\mu = \rr\to\qq$, $N(\qq) = \rr?x\seq B'$ for some $x$ and $B'$, $N'(\rr) = B$, $N'(\qq) = B'$, and $\sigma' = \upd{\sigma}{\qq.x}{v}$ where $\eval{e}{\sigma_\rr}{v}$.
  \item\label{it:inv-recv} $N(\rr) = \pp?x\seq B$; then the rule is \rn{N|Com}, $\mu = \pp\to\rr$, $N(\pp) = \rr!e\seq B'$ for some $e$ and $B'$, $N'(\rr) = B$, $N'(\pp) = B'$, and $\sigma' = \upd{\sigma}{\rr.x}{v}$ where $\eval{e}{\sigma_\pp}{v}$.
  \item\label{it:inv-cast} $N(\rr) = \sendbr{Q}{e}{B^1}{B^2}$; then the rule is \rn{N|Cast}, $\mu = \cast{\rr}{Q}$, $N(\qq) = \recvbr{\rr}{B_\qq^1}{B_\qq^2}$ for every $\qq\in Q$, $N'(\rr) = B^v$ and $N'(\qq) = B_\qq^v$ for every $\qq\in Q$, where $\eval{e}{\sigma_\rr}{v}$, and $\sigma' = \sigma$.
  \item\label{it:inv-brecv} $N(\rr) = \recvbr{\pp}{B_1}{B_2}$; then the rule is \rn{N|Cast}, $\mu = \cast{\pp}{Q}$ for some $Q$ with $\rr\in Q$, $N(\pp) = \sendbr{Q}{e}{B^1}{B^2}$ for some $e$, $B^1$, $B^2$, $N'(\rr) = B_v$ where $\eval{e}{\sigma_\pp}{v}$, and $\sigma' = \sigma$.
  \item\label{it:inv-call} $N(\rr) = X$; then the rule is \rn{N|Call}, $\mu = \tauat\rr$, $N'(\rr) = \ldef{\rr}{X}$, and $\sigma' = \sigma$.
  \end{enumerate}
\end{lemma}
\begin{proof}
  By inspection of Figure~\ref{fig:net-sem}.  In every rule, each process in
  $\pn(\mu)$ is displayed on the left-hand side with a behaviour of one of the
  six non-terminated shapes, so $N(\rr)\neq\nil$.  The six shapes are
  pairwise distinct, and each shape occurs in exactly one role of exactly one
  rule: assignment only as the subject of \rn{N|Local}, send and receive only
  as the sender and receiver of \rn{N|Com}, send-and-branch and
  receive-and-branch only as the sender and a receiver of \rn{N|Cast}, and
  call only as the subject of \rn{N|Call}.  Hence the shape of $N(\rr)$
  determines the rule and the role of $\rr$ in it, and the remaining claims
  in each item are read off from that rule.
\end{proof}

\begin{lemma}[Projection and process names]\label{lem:proj-prefix}
  \begin{enumerate}[label=(\roman*),nosep]
  \item\label{it:pp-act} If $\rr\notin\pn(\eta)$ then $\proj{\eta\seq C}{\rr} = \proj{C}{\rr}$.
  \item\label{it:pp-rt} For every $P\subseteq\Procs$, possibly empty,
    $\pn(\rt{P}{X}{C}) = P\cup\pn(C)$, and $\proj{\rt{P}{X}{C}}{\rr}$ is $X$
    if $\rr\in P$ and $\proj{C}{\rr}$ otherwise.
  \item\label{it:pp-nil} $\proj{C}{\rr} = \nil$ if and only if $\rr\notin\pn(C)$.
  \end{enumerate}
\end{lemma}
\begin{proof}
  \ref{it:pp-act} is immediate from Definition~\ref{def:epp}.
  \ref{it:pp-rt} is immediate from Definitions~\ref{def:pn}
  and~\ref{def:epp} when $P\neq\emptyset$; when $P=\emptyset$, the notation
  $\rt{P}{X}{C}$ stands for $C$, and $\rr\notin P$ for every $\rr$.
  \ref{it:pp-nil} is by induction on $C$.  For $\nil$ both sides hold.  For
  $\eta\seq C'$, if $\rr\in\pn(\eta)$ then $\proj{\eta\seq C'}{\rr}$ starts
  with an assignment, a send or a receive, and $\rr\in\pn(\eta\seq C')$;
  otherwise use \ref{it:pp-act} and the induction hypothesis.  For a
  conditional at $\pp$ with receivers $Q$, we have $\pn(C) = \{\pp\}\cup Q$,
  and the projection is a send-and-branch at $\pp$, a receive-and-branch at
  each $\rr\in Q$ and $\nil$ elsewhere.  For a call $X$, the claim is the
  definition of $\proj{X}{\rr}$.  For $\rt{P}{X}{C'}$, if $\rr\in P$ the
  projection is $X$ and $\rr\in\pn(C)$; otherwise the projection is
  $\proj{C'}{\rr}$ and $\rr\in\pn(C)$ iff $\rr\in\pn(C')$, and we use the
  induction hypothesis.
\end{proof}

\begin{lemma}[Projection invariance]\label{lem:proj-inv}
  If $\conf{C}{\sigma}\lto{\mu}\conf{C'}{\sigma'}$ and $\rr\notin\pn(\mu)$,
  then $\proj{C'}{\rr} = \proj{C}{\rr}$.
\end{lemma}
\begin{proof}
  By induction on the derivation of the transition.

  \case{\rn{C|Local}}
  $C = \pp.x := e\seq C'$ with $\rr\neq\pp$.  By
  Lemma~\ref{lem:proj-prefix}\ref{it:pp-act}, $\proj{C}{\rr} = \proj{C'}{\rr}$.

  \case{\rn{C|Com}}
  $C = \pp.e\to\qq.x\seq C'$ with $\rr\notin\{\pp,\qq\}$.  As in the previous
  case.

  \case{\rn{C|Cast}}
  $C = \ifcond{\pp.e}{C_1}{C_2}$, $C' = C_v$ and
  $\rr\notin\{\pp\}\cup Q = \pn(C)$.  Then $\proj{C}{\rr} = \nil$ by
  definition, and since $\rr\notin\pn(C_v)$, also $\proj{C_v}{\rr} = \nil$ by
  Lemma~\ref{lem:proj-prefix}\ref{it:pp-nil}.

  \case{\rn{C|Delay}}
  $C = \eta\seq C_0$ and $C' = \eta\seq C_0'$ with
  $\conf{C_0}{\sigma}\lto{\mu}\conf{C_0'}{\sigma'}$.  The induction
  hypothesis gives $\proj{C_0'}{\rr} = \proj{C_0}{\rr}$.  If $\rr\notin\pn(\eta)$
  the claim follows from Lemma~\ref{lem:proj-prefix}\ref{it:pp-act}.  If
  $\rr\in\pn(\eta)$, then by Definition~\ref{def:epp} both $\proj{C}{\rr}$ and
  $\proj{C'}{\rr}$ consist of the same prefix (an assignment, a send or a
  receive, determined by $\eta$ and $\rr$) followed by $\proj{C_0}{\rr}$,
  respectively $\proj{C_0'}{\rr}$, which are equal.

  \case{\rn{C|Call}}
  $C = X$, $\mu = \tauat\pp$ and $C' = \rt{(\pn(X)\setminus\{\pp\})}{X}{\body{X}}$,
  with $\rr\neq\pp$.  If $\rr\in\pn(X)$, then $\rr\in\pn(X)\setminus\{\pp\}$
  and both projections are $X$, by Definition~\ref{def:epp} and
  Lemma~\ref{lem:proj-prefix}\ref{it:pp-rt}.  Otherwise $\proj{C}{\rr} = \nil$,
  and $\proj{C'}{\rr} = \proj{\body{X}}{\rr} = \nil$ by
  Lemma~\ref{lem:proj-prefix}\ref{it:pp-rt} and~\ref{it:pp-nil}, since
  $\rr\notin\pn(X) = \pn(\body{X})$ by Lemma~\ref{lem:pn-body}.

  \case{\rn{C|Enter}}
  $C = \rt{P}{X}{C_0}$, $\mu = \tauat\pp$ and $C' = \rt{(P\setminus\{\pp\})}{X}{C_0}$,
  with $\rr\neq\pp$.  Then $\rr\in P$ iff $\rr\in P\setminus\{\pp\}$, so by
  Lemma~\ref{lem:proj-prefix}\ref{it:pp-rt} both projections are $X$ or both
  are $\proj{C_0}{\rr}$.

  \case{\rn{C|Inside}}
  $C = \rt{P}{X}{C_0}$ and $C' = \rt{P}{X}{C_0'}$ with
  $\conf{C_0}{\sigma}\lto{\mu}\conf{C_0'}{\sigma'}$.  If $\rr\in P$ both
  projections are $X$.  Otherwise they are $\proj{C_0}{\rr}$ and
  $\proj{C_0'}{\rr}$, which are equal by the induction hypothesis.
\end{proof}

\subsection{Well-formedness}
\label{sec:wf}

A process in $P$ that has not yet entered a call $\rt{P}{X}{C}$ is, in the
projection, still at its call to $X$; when it enters, it will continue as its
local body $\proj{\body{X}}{\rr}$.  For the choreography and the network to
stay in step, the body's current state $C$ must still expect exactly that
behaviour from $\rr$.

\begin{definition}[Well-formedness]\label{def:wf}
  The set of \emph{well-formed} choreographies is the least set such that
  $\nil$ and every call $X$ are well-formed; $\eta\seq C$ is well-formed if
  $C$ is; $\ifcond{\pp.e}{C_1}{C_2}$ is well-formed if $C_1$ and $C_2$ are; and
  $\rt{P}{X}{C}$ is well-formed if $C$ is well-formed and
  \[
    \proj{C}{\rr} = \proj{\body{X}}{\rr}\qquad\text{for every } \rr\in P.
  \]
  A configuration $\conf{C}{\sigma}$ is well-formed if $C$ is.
\end{definition}

Every source choreography is well-formed, since the only constraint concerns
runtime terms.  The last clause also holds for $P=\emptyset$ under our
convention: $\rt{\emptyset}{X}{C}$, which is $C$, is well-formed if and only
if $C$ is.

\begin{example}\label{ex:not-wf}
  Without well-formedness the correctness results fail.  Let
  $X = \pp.e\to\qq.x\seq X$ and $C = \rt{\{\qq\}}{X}{\nil}$, which is not
  well-formed because $\proj{\nil}{\qq} = \nil \neq \pp?x\seq X = \proj{\body{X}}{\qq}$.
  By \rn{C|Enter}, $\conf{C}{\sigma}\lto{\tauat\qq}\conf{\nil}{\sigma}$, which
  is terminated.  The projection $\projn{C}$ maps $\qq$ to $X$ and every
  other process to $\nil$; its only step unfolds $\qq$'s call to
  $\pp?x\seq X$, after which $\qq$ waits forever for a message from $\pp$.
\end{example}

\begin{lemma}[Preservation of well-formedness]\label{lem:wf-pres}
  If $C$ is well-formed and $\conf{C}{\sigma}\lto{\mu}\conf{C'}{\sigma'}$,
  then $C'$ is well-formed.
\end{lemma}
\begin{proof}
  By induction on the derivation of the transition.  For \rn{C|Local} and
  \rn{C|Com}, $C'$ is a subterm of $C$, and for \rn{C|Cast} it is a branch of
  $C$; subterms of well-formed choreographies are well-formed.  For
  \rn{C|Delay}, $C = \eta\seq C_0$ and $C' = \eta\seq C_0'$, and $C_0'$ is
  well-formed by the induction hypothesis.

  For \rn{C|Call}, $C' = \rt{(\pn(X)\setminus\{\pp\})}{X}{\body{X}}$.  The
  body $\body{X}$ is a source choreography, hence well-formed, and the
  required equation $\proj{\body{X}}{\rr} = \proj{\body{X}}{\rr}$ is trivial.

  For \rn{C|Enter}, $C = \rt{P}{X}{C_0}$ and
  $C' = \rt{(P\setminus\{\pp\})}{X}{C_0}$.  Since $C$ is well-formed, $C_0$ is
  well-formed and the required equation holds for every $\rr\in P$, so in
  particular for every $\rr\in P\setminus\{\pp\}$.

  For \rn{C|Inside}, $C = \rt{P}{X}{C_0}$ and $C' = \rt{P}{X}{C_0'}$ with
  $\conf{C_0}{\sigma}\lto{\mu}\conf{C_0'}{\sigma'}$ and
  $P\cap\pn(\mu)=\emptyset$.  The induction hypothesis gives that $C_0'$ is
  well-formed.  For $\rr\in P$ we have $\rr\notin\pn(\mu)$, so
  Lemma~\ref{lem:proj-inv} and the well-formedness of $C$ give
  $\proj{C_0'}{\rr} = \proj{C_0}{\rr} = \proj{\body{X}}{\rr}$.
\end{proof}

\subsection{Completeness}

\begin{theorem}[Completeness]\label{thm:completeness}
  If $C$ is well-formed and $\conf{C}{\sigma}\lto{\mu}\conf{C'}{\sigma'}$, then
  \[
    \conf{\projn{C}}{\sigma}\lto{\mu}\conf{\projn{C'}}{\sigma'}.
  \]
\end{theorem}
\begin{proof}
  By induction on the derivation of the transition.  Subterms of a
  well-formed choreography are well-formed, so the induction hypothesis
  applies to the premise of \rn{C|Delay} and of \rn{C|Inside}.

  \case{\rn{C|Local}}
  $C = \pp.x := e\seq C'$, $\mu = \tauat\pp$, $\eval{e}{\sigma_\pp}{v}$ and
  $\sigma' = \upd{\sigma}{\pp.x}{v}$.  We have $\proj{C}{\pp} = x := e\seq\proj{C'}{\pp}$
  and $\proj{C}{\rr} = \proj{C'}{\rr}$ for $\rr\neq\pp$.  Rule \rn{N|Local}
  gives $\conf{\projn{C}}{\sigma}\lto{\tauat\pp}\conf{N'}{\sigma'}$ with
  $N'(\pp) = \proj{C'}{\pp}$ and $N'(\rr) = \proj{C}{\rr} = \proj{C'}{\rr}$
  otherwise, so $N' = \projn{C'}$.

  \case{\rn{C|Com}}
  $C = \pp.e\to\qq.x\seq C'$, $\mu = \pp\to\qq$, $\eval{e}{\sigma_\pp}{v}$ and
  $\sigma' = \upd{\sigma}{\qq.x}{v}$.  We have
  $\proj{C}{\pp} = \qq!e\seq\proj{C'}{\pp}$, $\proj{C}{\qq} = \pp?x\seq\proj{C'}{\qq}$
  and $\proj{C}{\rr} = \proj{C'}{\rr}$ otherwise.  Rule \rn{N|Com} gives a
  transition labelled $\pp\to\qq$ to $\conf{\projn{C'}}{\sigma'}$.

  \case{\rn{C|Cast}}
  $C = \ifcond{\pp.e}{C_1}{C_2}$, $C' = C_v$, $\mu = \cast{\pp}{Q}$ with
  $Q = (\pn(C_1)\cup\pn(C_2))\setminus\{\pp\}$, $\eval{e}{\sigma_\pp}{v}$ and
  $\sigma' = \sigma$.  By Definition~\ref{def:epp},
  $\proj{C}{\pp} = \sendbr{Q}{e}{\proj{C_1}{\pp}}{\proj{C_2}{\pp}}$,
  $\proj{C}{\qq} = \recvbr{\pp}{\proj{C_1}{\qq}}{\proj{C_2}{\qq}}$ for $\qq\in Q$,
  and $\proj{C}{\rr} = \nil$ for $\rr\notin\{\pp\}\cup Q$.  Rule \rn{N|Cast}
  applies and gives $\conf{\projn{C}}{\sigma}\lto{\cast{\pp}{Q}}\conf{N'}{\sigma}$
  with $N'(\pp) = \proj{C_v}{\pp}$, $N'(\qq) = \proj{C_v}{\qq}$ for $\qq\in Q$
  and $N'(\rr) = \nil$ for $\rr\notin\{\pp\}\cup Q$.  For the last group,
  $\rr\notin\pn(C_v)$, so $\proj{C_v}{\rr} = \nil = N'(\rr)$ by
  Lemma~\ref{lem:proj-prefix}\ref{it:pp-nil}.  Hence $N' = \projn{C_v} = \projn{C'}$.

  \case{\rn{C|Delay}}
  $C = \eta\seq C_0$ and $C' = \eta\seq C_0'$ with
  $\conf{C_0}{\sigma}\lto{\mu}\conf{C_0'}{\sigma'}$ and
  $\pn(\eta)\cap\pn(\mu) = \emptyset$.  The induction hypothesis gives
  $\conf{\projn{C_0}}{\sigma}\lto{\mu}\conf{\projn{C_0'}}{\sigma'}$.  By
  Lemma~\ref{lem:proj-prefix}\ref{it:pp-act}, $\projn{C}$ agrees with
  $\projn{C_0}$ on $\pn(\mu)$, since $\pn(\mu)$ is disjoint from $\pn(\eta)$.
  By Lemma~\ref{lem:frame}\ref{it:frame-frame},
  $\conf{\projn{C}}{\sigma}\lto{\mu}\conf{N'}{\sigma'}$ where
  $N'(\rr) = \proj{C_0'}{\rr}$ for $\rr\in\pn(\mu)$ and $N'(\rr) = \proj{C}{\rr}$
  otherwise.  It remains to show $N' = \projn{C'}$.  For $\rr\in\pn(\mu)$ we
  have $\rr\notin\pn(\eta)$, so $\proj{C'}{\rr} = \proj{C_0'}{\rr} = N'(\rr)$.
  For $\rr\notin\pn(\mu)$, Lemma~\ref{lem:proj-inv} gives
  $\proj{C_0'}{\rr} = \proj{C_0}{\rr}$, and therefore
  $\proj{\eta\seq C_0'}{\rr} = \proj{\eta\seq C_0}{\rr} = N'(\rr)$, since both
  sides either add the same prefix to equal projections (if $\rr\in\pn(\eta)$)
  or are the equal projections themselves.

  \case{\rn{C|Call}}
  $C = X$, $\mu = \tauat\pp$ with $\pp\in\pn(X)$,
  $C' = \rt{(\pn(X)\setminus\{\pp\})}{X}{\body{X}}$ and $\sigma' = \sigma$.
  Since $\proj{C}{\pp} = X$, rule \rn{N|Call} gives
  $\conf{\projn{C}}{\sigma}\lto{\tauat\pp}\conf{N'}{\sigma}$ with
  $N'(\pp) = \ldef{\pp}{X} = \proj{\body{X}}{\pp}$ and $N'(\rr) = \proj{C}{\rr}$
  for $\rr\neq\pp$.  Since $\pp\notin\pn(X)\setminus\{\pp\}$,
  Lemma~\ref{lem:proj-prefix}\ref{it:pp-rt} gives
  $\proj{C'}{\pp} = \proj{\body{X}}{\pp} = N'(\pp)$, and for $\rr\neq\pp$
  Lemma~\ref{lem:proj-inv} gives $\proj{C'}{\rr} = \proj{C}{\rr} = N'(\rr)$.
  Hence $N' = \projn{C'}$.

  \case{\rn{C|Enter}}
  $C = \rt{P}{X}{C_0}$, $\mu = \tauat\pp$ with $\pp\in P$,
  $C' = \rt{(P\setminus\{\pp\})}{X}{C_0}$ and $\sigma' = \sigma$.  Since
  $\proj{C}{\pp} = X$, rule \rn{N|Call} gives
  $\conf{\projn{C}}{\sigma}\lto{\tauat\pp}\conf{N'}{\sigma}$ with
  $N'(\pp) = \proj{\body{X}}{\pp}$ and $N'(\rr) = \proj{C}{\rr}$ for
  $\rr\neq\pp$.  By Lemma~\ref{lem:proj-prefix}\ref{it:pp-rt} and the
  well-formedness of $C$, $\proj{C'}{\pp} = \proj{C_0}{\pp} = \proj{\body{X}}{\pp} = N'(\pp)$.
  For $\rr\neq\pp$, Lemma~\ref{lem:proj-inv} gives
  $\proj{C'}{\rr} = \proj{C}{\rr} = N'(\rr)$.  Hence $N' = \projn{C'}$.

  \case{\rn{C|Inside}}
  $C = \rt{P}{X}{C_0}$ and $C' = \rt{P}{X}{C_0'}$ with
  $\conf{C_0}{\sigma}\lto{\mu}\conf{C_0'}{\sigma'}$ and
  $P\cap\pn(\mu) = \emptyset$.  The induction hypothesis gives
  $\conf{\projn{C_0}}{\sigma}\lto{\mu}\conf{\projn{C_0'}}{\sigma'}$.  By
  Definition~\ref{def:epp}, $\projn{C}$ agrees with $\projn{C_0}$ on
  $\pn(\mu)$, since $\pn(\mu)$ is disjoint from $P$.  By
  Lemma~\ref{lem:frame}\ref{it:frame-frame},
  $\conf{\projn{C}}{\sigma}\lto{\mu}\conf{N'}{\sigma'}$ where
  $N'(\rr) = \proj{C_0'}{\rr}$ for $\rr\in\pn(\mu)$ and
  $N'(\rr) = \proj{C}{\rr}$ otherwise.  For $\rr\in\pn(\mu)$ we have
  $\rr\notin P$, so $\proj{C'}{\rr} = \proj{C_0'}{\rr} = N'(\rr)$.  For
  $\rr\notin\pn(\mu)$, Lemma~\ref{lem:proj-inv} gives
  $\proj{C'}{\rr} = \proj{C}{\rr} = N'(\rr)$.  Hence $N' = \projn{C'}$.
\end{proof}

\subsection{Soundness}

\begin{theorem}[Soundness]\label{thm:soundness}
  If $C$ is well-formed and $\conf{\projn{C}}{\sigma}\lto{\mu}\conf{N'}{\sigma'}$,
  then there is a choreography $C'$ such that
  \[
    \conf{C}{\sigma}\lto{\mu}\conf{C'}{\sigma'}
    \qquad\text{and}\qquad
    N' = \projn{C'}.
  \]
\end{theorem}
\begin{proof}
  By structural induction on $C$.  Throughout, $N$ denotes $\projn{C}$.
  Subterms of a well-formed choreography are well-formed, so the induction
  hypothesis applies to them.

  \case{$C = \nil$}
  Every process of $N$ is $\nil$, so by Lemma~\ref{lem:inversion} no
  transition is possible and there is nothing to prove.

  \case{$C = \eta\seq C_0$}
  We distinguish whether the step involves a process of $\eta$.

  \subcase{$\pn(\eta)\cap\pn(\mu) \neq \emptyset$}
  First let $\eta = \pp.x := e$, so that $\pp\in\pn(\mu)$ and
  $N(\pp) = x := e\seq\proj{C_0}{\pp}$.  By
  Lemma~\ref{lem:inversion}\ref{it:inv-local}, $\mu = \tauat\pp$,
  $\sigma' = \upd{\sigma}{\pp.x}{v}$ with $\eval{e}{\sigma_\pp}{v}$, and
  $N'(\pp) = \proj{C_0}{\pp}$; by Lemma~\ref{lem:frame}\ref{it:frame-loc},
  $N'(\rr) = N(\rr) = \proj{C_0}{\rr}$ for $\rr\neq\pp$
  (Lemma~\ref{lem:proj-prefix}\ref{it:pp-act}).  Hence $N' = \projn{C_0}$,
  and \rn{C|Local} gives $\conf{C}{\sigma}\lto{\tauat\pp}\conf{C_0}{\sigma'}$;
  take $C' = C_0$.

  Now let $\eta = \pp.e\to\qq.x$, so that $N(\pp) = \qq!e\seq\proj{C_0}{\pp}$
  and $N(\qq) = \pp?x\seq\proj{C_0}{\qq}$, and $\pp\in\pn(\mu)$ or
  $\qq\in\pn(\mu)$.  In the first case Lemma~\ref{lem:inversion}\ref{it:inv-send}
  applies to $\pp$; in the second case Lemma~\ref{lem:inversion}\ref{it:inv-recv}
  applies to $\qq$.  Either way we obtain $\mu = \pp\to\qq$,
  $\sigma' = \upd{\sigma}{\qq.x}{v}$ with $\eval{e}{\sigma_\pp}{v}$,
  $N'(\pp) = \proj{C_0}{\pp}$ and $N'(\qq) = \proj{C_0}{\qq}$, and by
  Lemma~\ref{lem:frame}\ref{it:frame-loc} and
  Lemma~\ref{lem:proj-prefix}\ref{it:pp-act}, $N'(\rr) = \proj{C_0}{\rr}$ for
  all other $\rr$.  Hence $N' = \projn{C_0}$, and \rn{C|Com} gives
  $\conf{C}{\sigma}\lto{\pp\to\qq}\conf{C_0}{\sigma'}$; take $C' = C_0$.

  \subcase{$\pn(\eta)\cap\pn(\mu) = \emptyset$}
  By Lemma~\ref{lem:proj-prefix}\ref{it:pp-act}, $\projn{C_0}$ agrees with
  $N$ on $\pn(\mu)$.  By Lemma~\ref{lem:frame}\ref{it:frame-frame},
  $\conf{\projn{C_0}}{\sigma}\lto{\mu}\conf{N''}{\sigma'}$ where $N''$ agrees
  with $N'$ on $\pn(\mu)$ and with $\projn{C_0}$ elsewhere.  The induction
  hypothesis yields $C_0'$ with $\conf{C_0}{\sigma}\lto{\mu}\conf{C_0'}{\sigma'}$
  and $\projn{C_0'} = N''$.  Rule \rn{C|Delay} gives
  $\conf{\eta\seq C_0}{\sigma}\lto{\mu}\conf{\eta\seq C_0'}{\sigma'}$; take
  $C' = \eta\seq C_0'$.  It remains to show $\projn{C'} = N'$.  For
  $\rr\in\pn(\mu)$ we have $\rr\notin\pn(\eta)$, so
  $\proj{C'}{\rr} = \proj{C_0'}{\rr} = N''(\rr) = N'(\rr)$.  For
  $\rr\notin\pn(\mu)$, Lemma~\ref{lem:frame}\ref{it:frame-loc} gives
  $N'(\rr) = N(\rr) = \proj{\eta\seq C_0}{\rr}$, and Lemma~\ref{lem:proj-inv}
  gives $\proj{C_0'}{\rr} = \proj{C_0}{\rr}$, hence
  $\proj{\eta\seq C_0'}{\rr} = \proj{\eta\seq C_0}{\rr} = N'(\rr)$.

  \case{$C = \ifcond{\pp.e}{C_1}{C_2}$}
  Let $Q = (\pn(C_1)\cup\pn(C_2))\setminus\{\pp\}$.  By
  Definition~\ref{def:epp}, $N(\pp) = \sendbr{Q}{e}{\proj{C_1}{\pp}}{\proj{C_2}{\pp}}$,
  $N(\qq) = \recvbr{\pp}{\proj{C_1}{\qq}}{\proj{C_2}{\qq}}$ for $\qq\in Q$, and
  $N(\rr) = \nil$ otherwise.  By Lemma~\ref{lem:inversion} every process in
  $\pn(\mu)$ has a non-terminated behaviour, so $\pn(\mu)\subseteq\{\pp\}\cup Q$.
  We claim $\pp\in\pn(\mu)$.  Otherwise $\pn(\mu)$ contains some $\qq\in Q$,
  whose behaviour is a receive-and-branch from $\pp$, and
  Lemma~\ref{lem:inversion}\ref{it:inv-brecv} gives $\mu = \cast{\pp}{Q'}$ for
  some $Q'$, contradicting $\pp\notin\pn(\mu)$.  So $\pp\in\pn(\mu)$ and, since
  $N(\pp)$ is a send-and-branch, Lemma~\ref{lem:inversion}\ref{it:inv-cast}
  gives $\mu = \cast{\pp}{Q}$, $\sigma' = \sigma$, $N'(\pp) = \proj{C_v}{\pp}$
  and $N'(\qq) = \proj{C_v}{\qq}$ for $\qq\in Q$, where $\eval{e}{\sigma_\pp}{v}$;
  by Lemma~\ref{lem:frame}\ref{it:frame-loc}, $N'(\rr) = \nil$ for
  $\rr\notin\{\pp\}\cup Q$.  Rule \rn{C|Cast} gives
  $\conf{C}{\sigma}\lto{\cast{\pp}{Q}}\conf{C_v}{\sigma}$; take $C' = C_v$.
  Finally $\projn{C_v} = N'$: at $\pp$ and at every $\qq\in Q$ both are
  $\proj{C_v}{\cdot}$ by the above, and at $\rr\notin\{\pp\}\cup Q$ we have
  $\rr\notin\pn(C_v)$, so $\proj{C_v}{\rr} = \nil = N'(\rr)$ by
  Lemma~\ref{lem:proj-prefix}\ref{it:pp-nil}.

  \case{$C = X$}
  Let $\rr\in\pn(\mu)$; every label involves at least one process.  By
  Lemma~\ref{lem:inversion}, $N(\rr)\neq\nil$, so $\rr\in\pn(X)$ and
  $N(\rr) = X$.  By Lemma~\ref{lem:inversion}\ref{it:inv-call},
  $\mu = \tauat\rr$, $\sigma' = \sigma$ and
  $N'(\rr) = \ldef{\rr}{X} = \proj{\body{X}}{\rr}$, and by
  Lemma~\ref{lem:frame}\ref{it:frame-loc}, $N'(\qq) = N(\qq)$ for
  $\qq\neq\rr$.  Rule \rn{C|Call} gives
  $\conf{X}{\sigma}\lto{\tauat\rr}\conf{C'}{\sigma}$ with
  $C' = \rt{(\pn(X)\setminus\{\rr\})}{X}{\body{X}}$.  By
  Lemma~\ref{lem:proj-prefix}\ref{it:pp-rt},
  $\proj{C'}{\rr} = \proj{\body{X}}{\rr} = N'(\rr)$, and for $\qq\neq\rr$
  Lemma~\ref{lem:proj-inv} gives $\proj{C'}{\qq} = \proj{C}{\qq} = N(\qq) = N'(\qq)$.
  Hence $N' = \projn{C'}$.

  \case{$C = \rt{P}{X}{C_0}$}
  By Definition~\ref{def:epp}, $N(\rr) = X$ for $\rr\in P$ and
  $N(\rr) = \proj{C_0}{\rr}$ otherwise.  We distinguish whether the step
  involves a process that has not entered the call.

  \subcase{$P\cap\pn(\mu)\neq\emptyset$}
  Let $\rr\in P\cap\pn(\mu)$.  Since $N(\rr) = X$,
  Lemma~\ref{lem:inversion}\ref{it:inv-call} gives $\mu = \tauat\rr$,
  $\sigma' = \sigma$ and $N'(\rr) = \proj{\body{X}}{\rr}$, and by
  Lemma~\ref{lem:frame}\ref{it:frame-loc}, $N'(\qq) = N(\qq)$ for
  $\qq\neq\rr$.  Rule \rn{C|Enter} gives
  $\conf{C}{\sigma}\lto{\tauat\rr}\conf{C'}{\sigma}$ with
  $C' = \rt{(P\setminus\{\rr\})}{X}{C_0}$.  By
  Lemma~\ref{lem:proj-prefix}\ref{it:pp-rt} and the well-formedness of $C$,
  $\proj{C'}{\rr} = \proj{C_0}{\rr} = \proj{\body{X}}{\rr} = N'(\rr)$, and
  for $\qq\neq\rr$ Lemma~\ref{lem:proj-inv} gives
  $\proj{C'}{\qq} = \proj{C}{\qq} = N(\qq) = N'(\qq)$.  Hence
  $N' = \projn{C'}$.

  \subcase{$P\cap\pn(\mu) = \emptyset$}
  Then $\projn{C_0}$ agrees with $N$ on $\pn(\mu)$.  By
  Lemma~\ref{lem:frame}\ref{it:frame-frame},
  $\conf{\projn{C_0}}{\sigma}\lto{\mu}\conf{N''}{\sigma'}$ where $N''$
  agrees with $N'$ on $\pn(\mu)$ and with $\projn{C_0}$ elsewhere.  The
  induction hypothesis yields $C_0'$ with
  $\conf{C_0}{\sigma}\lto{\mu}\conf{C_0'}{\sigma'}$ and $\projn{C_0'} = N''$.
  Rule \rn{C|Inside} gives
  $\conf{C}{\sigma}\lto{\mu}\conf{\rt{P}{X}{C_0'}}{\sigma'}$; take
  $C' = \rt{P}{X}{C_0'}$.  For $\rr\in\pn(\mu)$ we have $\rr\notin P$, so
  $\proj{C'}{\rr} = \proj{C_0'}{\rr} = N''(\rr) = N'(\rr)$.  For
  $\rr\notin\pn(\mu)$, Lemma~\ref{lem:proj-inv} and
  Lemma~\ref{lem:frame}\ref{it:frame-loc} give
  $\proj{C'}{\rr} = \proj{C}{\rr} = N(\rr) = N'(\rr)$.  Hence
  $N' = \projn{C'}$.
\end{proof}

\subsection{Consequences}

Together, Theorems~\ref{thm:completeness} and~\ref{thm:soundness} and
Lemma~\ref{lem:wf-pres} say that projection is a functional bisimulation on
well-formed configurations.

\begin{corollary}[Bisimulation]\label{cor:bisim}
  Let $\Bis$ be the relation
  \[
    \Bis = \{\,(\conf{C}{\sigma},\, \conf{\projn{C}}{\sigma}) \mid C \text{ a well-formed choreography and } \sigma \text{ a store}\,\}.
  \]
  Then $\Bis$ is a strong labelled bisimulation between the transition
  systems of Definitions~\ref{def:chor-sem} and~\ref{def:net-sem}, the latter
  taken with respect to $\projn{\Defs}$:
  \begin{enumerate}[label=(\roman*),nosep]
  \item if $(\conf{C}{\sigma}, \conf{N}{\sigma})\in\Bis$ and
    $\conf{C}{\sigma}\lto{\mu}\conf{C'}{\sigma'}$, then
    $\conf{N}{\sigma}\lto{\mu}\conf{N'}{\sigma'}$ for some $N'$ with
    $(\conf{C'}{\sigma'}, \conf{N'}{\sigma'})\in\Bis$;
  \item if $(\conf{C}{\sigma}, \conf{N}{\sigma})\in\Bis$ and
    $\conf{N}{\sigma}\lto{\mu}\conf{N'}{\sigma'}$, then
    $\conf{C}{\sigma}\lto{\mu}\conf{C'}{\sigma'}$ for some $C'$ with
    $(\conf{C'}{\sigma'}, \conf{N'}{\sigma'})\in\Bis$.
  \end{enumerate}
\end{corollary}
\begin{proof}
  (i) is Theorem~\ref{thm:completeness} with $N' = \projn{C'}$, and (ii) is
  Theorem~\ref{thm:soundness}.  In both cases $C'$ is well-formed by
  Lemma~\ref{lem:wf-pres}.
\end{proof}

Unlike the corresponding result for CC, the network side of the bisimulation
is the projection itself and not a network related to it by pruning: since
every participant of a conditional learns the outcome in the very step that
resolves it, the projection never contains code for a branch that has
already been ruled out.  Procedure calls are matched step for step as well,
rather than up to the unfolding of calls, because entering a call is a
labelled step of both systems.  The only hypothesis is well-formedness, which
every source choreography satisfies.

\begin{lemma}[Progress]\label{lem:progress}
  If $\pn(C)\neq\emptyset$ then $\conf{C}{\sigma}\lto{\mu}\conf{C'}{\sigma'}$
  for some $\mu$, $C'$ and $\sigma'$.
\end{lemma}
\begin{proof}
  By cases on $C$, using Assumption~\ref{ass:eval}.  The case $C = \nil$ is
  excluded since $\pn(\nil) = \emptyset$.  If $C = \pp.x := e\seq C_0$
  or $C = \pp.e\to\qq.x\seq C_0$, then $e$ evaluates to some $v$ in
  $\sigma_\pp$ and \rn{C|Local}, respectively \rn{C|Com}, applies.  If
  $C = \ifcond{\pp.e}{C_1}{C_2}$, then $e$ evaluates to a boolean $v$ in
  $\sigma_\pp$ and \rn{C|Cast} applies.  If $C = X$, then
  $\pn(X) = \pn(C)\neq\emptyset$ and \rn{C|Call} applies for any
  $\pp\in\pn(X)$.  If $C = \rt{P}{X}{C_0}$, then $P\neq\emptyset$ and
  \rn{C|Enter} applies for any $\pp\in P$.
\end{proof}

\begin{corollary}[Deadlock-freedom by construction]\label{cor:df}
  For every well-formed choreography $C$ and store $\sigma$, the
  configuration $\conf{\projn{C}}{\sigma}$ has no transition if and only if
  $\projn{C}$ is terminated, if and only if $\pn(C) = \emptyset$.
  Consequently, every network configuration reachable from the projection of
  a well-formed choreography, and in particular of a source choreography,
  either has a transition or is terminated.
\end{corollary}
\begin{proof}
  By Lemma~\ref{lem:proj-prefix}\ref{it:pp-nil}, $\projn{C}$ is terminated
  if and only if $\pn(C) = \emptyset$.  If $\projn{C}$ is terminated then, by
  Lemma~\ref{lem:inversion}, it has no transition.  Otherwise
  $\pn(C)\neq\emptyset$, so by Lemma~\ref{lem:progress} $\conf{C}{\sigma}$
  has a transition, and by Theorem~\ref{thm:completeness} so does
  $\conf{\projn{C}}{\sigma}$.  This establishes the two equivalences.  For
  the last claim, by Corollary~\ref{cor:bisim} every configuration reachable
  from $\conf{\projn{C}}{\sigma}$ is of the form $\conf{\projn{C'}}{\sigma'}$
  for some well-formed choreography $C'$, to which the equivalences apply.
\end{proof}

\begin{remark}[Procedures without participants]\label{rem:empty}
  The condition $\pn(C) = \emptyset$ does not imply $C = \nil$.  A procedure
  whose body, and the bodies of all procedures reachable from it, contain no
  actions and no conditionals, such as $X = X$, has no participants.  A call
  to it has no transitions and projects to $\nil$ at every process; it
  describes a computation in which no process ever does anything.  If every
  procedure has at least one participant, then $\pn(C) = \emptyset$ if and
  only if $C = \nil$, and Corollary~\ref{cor:df} takes the same form as for
  the language without procedures.
\end{remark}

\end{document}
